% Generated draft: requires mathematical and source review.
\documentclass[11pt]{article}
\usepackage[margin=25mm]{geometry}
\usepackage{fontspec}
\setmainfont{Latin Modern Roman}
\usepackage{amsmath,amssymb,mathtools,unicode-math}
\setmathfont{Latin Modern Math}
\usepackage{enumitem,needspace,xurl,hyperref}
\setlength{\parindent}{0pt}
\setlength{\parskip}{5pt}
\setlength{\emergencystretch}{3em}
\newcommand{\N}{\mathbb{N}}
\newcommand{\Z}{\mathbb{Z}}
\newcommand{\Q}{\mathbb{Q}}
\newcommand{\R}{\mathbb{R}}
\newcommand{\C}{\mathbb{C}}
\newcommand{\F}{\mathbb{F}}
\newcommand{\A}{\mathbb{A}}
\newcommand{\PP}{\mathbb{P}}
\newcommand{\E}{\mathbb{E}}
\newcommand{\eps}{\varepsilon}
\newcommand{\dd}{\,\mathrm{d}}
\begin{document}
\section*{UnsolvedMath ID: 10600039}
\textbf{Source problem code:} catalogue entry
\subsection{Definitions and mathematical statement}
A chord diagram consists of a circle with finitely many distinct marked points paired into chords, without crossing signs or arrow directions in the free-knot setting. Two chords are linked when their endpoints alternate around the circle. The interlacement graph has one vertex per chord and an edge precisely for a linked pair. A free knot is represented by such a diagram modulo the free-knot Reidemeister equivalence; the trivial free knot is represented by the chordless circle.

Let $D$ be the source's heptagon Gauss diagram, with seven chords $c_0,\ldots,c_6$ whose linking satisfies
\[
 c_i\text{ is linked with }c_j
 \quad\Longleftrightarrow\quad j-i\equiv 1\text{ or }-1\pmod 7.
\]
Thus its interlacement graph is the cycle $C_7$. Problem-list item 39 asks whether the free knot represented by $D$ is trivial, that is, whether there is a finite sequence of allowed free Reidemeister moves taking $D$ to the chordless circle.

The supplied description specifies the interlacement graph but gives no drawing or cyclic endpoint word for the particular diagram. It must therefore be checked that this description identifies the intended free-knot class, or else the intended realization must be supplied. The excerpt does not explicitly ask a universal question about every diagram realizing $C_7$.
\subsection{Short English statement}
Decide whether the seven-chord free knot whose chord-linking graph is a heptagon can be reduced to the trivial free knot by free Reidemeister moves.
\subsection{Sources}
\begin{itemize}
\item Local registry: \url{data/unsolvedmath/problems.json}, permanent UnsolvedMath ID $10600039$, exact source code \texttt{catalogue entry}.
\item \emph{Unsolved Problems in Virtual Knot Theory and Combinatorial Knot Theory}, revised survey, problem-list item 39: \url{https://arxiv.org/abs/1409.2823}. Supplied access date: 2026-07-29; extraction: source-tex; the source list has 60 top-level items.
\item No separate primary citation, exact diagram, or endpoint word is supplied. The registry reports searches for ``free knot heptagon'' without a resolution; this is a supplied search assessment.
\end{itemize}
\end{document}
